% APPENDIX B: bảng/hình bổ sung, mã nguồn và dữ liệu (liên kết kho lưu trữ, bản phát hành).
\chapter{SUPPLEMENTARY MATERIALS}\label{app:supplementary}

\section{Source Code and Data Availability}
The research code, the application and this document are kept in one private version-controlled
repository held by the author; access for examination is arranged through the supervisor. The
repository contains the analysis scripts, the locked protocols, the generated result files and the
asset generators that produce every number in this document, so each figure can be traced to the
file it was read from.

Large artefacts are deliberately not in version control: model checkpoints, cached representations,
out-of-fold prediction files and the research data bundles. They are reproducible from the recorded
protocols and runner configurations, and each is pinned by size and SHA-256 in the manifests that
accompany the results. Primary data are not redistributed here; they are obtained from the original
public sources cited in Chapter~\ref{chap:background}, under the terms those sources set. Sources
whose terms do not permit redistribution are used only where those terms allow and are excluded from
any packaged bundle.

\section{Controls Moved from the Results Chapter}
The three studies below were completed under criteria fixed in advance and are reported here in full.
They are not part of the argument of Chapter~\ref{chap:evaluation}: two produced results that did not
support the intervention they tested, and the third establishes only how little this benchmark
overlaps a curated motif resource. Their numbers, denominators and intervals are unchanged from the
original analysis.

\section{Sensitivity to the Clustering Granularity}\label{sec:res-recluster}
Because cov50 does not change the partition on this peptide set, a grid of seven partitions of the 176 peptides
(the \nNPep{} benchmark peptides and the 85 added ones) was fixed before running: grouping by UniProt accession
(no clustering), the reference MMseqs2 setting (30\% identity, 80\% coverage) and five looser MMseqs2 settings. The P0
comparison (ridge M3, relaxed threshold 4.0 against the benchmark training set, same cluster bootstrap) was repeated with the
folds rebuilt from each partition. The reference level reproduces the earlier scores (\nAucRidge{} and \nAucRidgeExt) and the five
folds of the earlier experiment. The criterion fixed in advance: the gain is ``robust to granularity'' if every
partition that differs from the reference and is still informative (fewer than half of the added cells are lost because they
share a group with a test peptide) has a bootstrap interval of the gain above zero and a positive gain in all three seeds.

\begin{table}[htbp]
  \centering
  \caption[Gain of the relaxed threshold 4.0 (ridge M3) under different protein partitions]{Gain of the relaxed threshold 4.0 (ridge M3) under different protein partitions. ``Groups'' are the groups of the
  91 benchmark peptides and of all 176 peptides; ``Shared'' is the number of added peptides that fall in a group of a
  benchmark peptide; ``Excl.'' is the mean percentage of the added cells that are removed from training because their
  group contains a test peptide.}\label{tab:recluster}
  % Generated by build_ch4_assets.py - do not edit.
\begingroup\footnotesize\setlength{\tabcolsep}{3pt}
\renewcommand{\arraystretch}{1.15}
\begin{tabularx}{\textwidth}{P{3.4cm}rrrP{2.9cm}>{\raggedright\arraybackslash}Xr}
\toprule
Partition (MMseqs2) & Groups base/all & Shared & Excl.\,\% & AUROC base $\to$ relaxed & Gain [95\% CI] & Seeds $>0$ \\
\midrule
Accession (no clustering) & 91/176 & 0 & 0.0 & 0.8036 $\to$ 0.8202 & +0.0166 [+0.0099, +0.0240] & yes \\
id 0.3, cov 0.8 (reference) & 89/172 & 2 & 0.5 & 0.8082 $\to$ 0.8208 & +0.0126 [+0.0061, +0.0190] & yes \\
id 0.3, cov 0.5 \newline {\scriptsize grouping as reference; folds assigned differently} & 89/172 & 2 & 0.5 & 0.8116 $\to$ 0.8212 & +0.0096 [+0.0032, +0.0159] & yes \\
id 0.2, cov 0.5 & 87/163 & 6 & 1.4 & 0.8051 $\to$ 0.8185 & +0.0134 [+0.0074, +0.0199] & yes \\
id 0.1, cov 0.5 \newline {\scriptsize identical to the row above} & 87/163 & 6 & 1.4 & 0.8051 $\to$ 0.8185 & +0.0134 [+0.0074, +0.0199] & yes \\
id 0.1, cov 0.3 & 71/133 & 22 & 5.2 & 0.8096 $\to$ 0.8209 & +0.0113 [+0.0037, +0.0184] & yes \\
id 0.0, cov 0.3 \newline {\scriptsize identical to the row above} & 71/133 & 22 & 5.2 & 0.8096 $\to$ 0.8209 & +0.0113 [+0.0037, +0.0184] & yes \\
\bottomrule
\end{tabularx}
\par\smallskip{\footnotesize Source: \nolinkurl{recluster-sensitivity-2026-10-22/results.json}.}
\endgroup

\end{table}

The seven levels give only \nRcNDistinct{} distinct partitions (Table~\ref{tab:recluster}): lowering the identity from 0.2 to 0.1
at coverage 0.5, and from 0.1 to 0.0 at coverage 0.3, changes no group, and 30\% identity with coverage 0.5 reproduces the
reference grouping (only the group labels, and hence the fold assignment, differ; the gain there is
\nRcFiftyD{} \nRcFiftyCI{}, the same value as the cov50 row of Table~\ref{tab:partcrobust}). Three partitions differ from the
reference and are informative: grouping by accession (\nRcAccD{}, \nRcAccCI{}), 20\% identity with coverage 0.5 (\nRcMidD{}, \nRcMidCI{}) and
10\% identity with coverage 0.3 (\nRcCoarseD{}, \nRcCoarseCI{}). All intervals are above zero and the gain is positive in all
three seeds, so under the fixed rule the gain is \emph{robust to granularity}. The point estimates are not monotone in the
coarseness (the largest is for grouping by accession), the intervals overlap, and no test of their differences was
planned; the absolute benchmark score also moves with the partition (\nAucRidge{} for the reference and \nRcFiftyBase{} for the same grouping with
other folds), so absolute values must not be compared across rows.

The test is weaker than its name suggests. The coarsest partition still has \nRcCoarseGroups{} groups for the 91 benchmark peptides; in
it, \nRcCoarseShared{} added peptides share a group with a benchmark peptide and on average \nRcCoarseExcl\% of the added cells are removed
from training, and one of the five test folds consists of a single group. Whether the gain would survive a still coarser
split, for example by protein family, is not known.


\section{Pre-Training on Tile Labels}\label{sec:res-pretrain}
The experiment of Voutsinos et al.\ \citep{voutsinos2026} provides a degradation score for \nPtNTiles{} tiles of 30 residues,
without substitutions. The question was whether initialising the gated network (without $\Delta$PAP) from a model
pre-trained to predict these scores lets it learn $\Delta$PSI better under the same label budget. The design was fixed in
advance. Pre-training treats each tile as 30 wild-type rows (the substituted-residue input is zero), predicts the tile
as the mean over positions and regresses the sign-reversed, standardised degron score. Two branches then share the same
architecture, the learning-rate grid ($2\times10^{-3}$ or $5\times10^{-4}$, chosen on the inner validation set) and the
inner validation split, and differ only in whether the network starts from the pre-trained state (\emph{pre}) or from
random initialisation (\emph{rand}). To prevent leakage, in each outer fold every gene that has a tile similar to a protein of
an outer-test cluster (MMseqs2 at 30\% identity with an alignment of at least 50 residues, a shared 12-mer, or the same gene
name) is removed entirely from the pre-training data \emph{before} any learning; this removed \nPtGenesExcl{} genes in the five folds.
The gate on the reproduction of earlier scores was passed (ridge \nAucRidge{} and \nAucRidgeExt; the gated network within \nPtGateGap{}
of the earlier values). Two comparisons were fixed, both with the rule ``bootstrap interval above zero and a positive
difference in every seed'': (A) \emph{pre} against \emph{rand}, and (B) \emph{pre} against ridge M3.

The result is negative (Tables~\ref{tab:p3auroc} and~\ref{tab:p3verdict}). In the main configuration the pre-trained network reaches
\nPtExtPre{}, against \nPtExtRand{} for random initialisation: difference A is \nPtExtA{} \nPtExtACI{}, an interval that contains zero,
so pre-training cannot be distinguished from random initialisation; difference B is \nPtExtB{} \nPtExtBCI{}, so the pre-trained network
is below ridge. At threshold 3.0 both differences are clearly negative (A: \nPtBaseA{} \nPtBaseACI{}; B: \nPtBaseB{} \nPtBaseBCI{}).
The difference is negative in all three seeds in every comparison. By the mapping fixed in advance, the conclusion is
``insufficient evidence for a benefit of pre-training''. The pre-training task itself was learned (Pearson correlation
\nPtPearson{} between prediction and score on held-out genes, range \nPtPearsonLo{}--\nPtPearsonHi{} over the fifteen runs), so learning this source endpoint does not by itself
establish useful transfer to $\Delta$PSI. This result rejects only this design, in which the
pre-training sees wild-type sequences only and the whole network is then fine-tuned; designs that pre-train with
substitutions, freeze part of this tile-label encoder or train both tasks jointly were not tried.
The later frozen-ESM and LoRA experiments use different pretrained representations and are reported separately.

\begin{table}[htbp]
  \centering
  \caption[Pre-registered comparisons for pre-training]{Pre-registered comparisons for pre-training. The main configuration is the relaxed threshold 4.0; the threshold 3.0 is
  reported with the same criteria and was not used to choose the configuration.}\label{tab:p3verdict}
  % Generated by build_ch4_assets.py - do not edit.
\begingroup\footnotesize
\renewcommand{\arraystretch}{1.15}
\begin{tabularx}{\textwidth}{P{1.6cm}P{3.0cm}P{3.5cm}P{3.5cm}r}
\toprule
Thr. & Comparison & Difference [95\% CI] & Per-seed difference & Met \\
\midrule
4.0 (main) & A: pre-trained $-$ random init. & $-$0.0077 [$-$0.0163, +0.0006] & $-$0.0180, $-$0.0259, $-$0.0121 & no \\
 & B: pre-trained $-$ ridge M3 & $-$0.0110 [$-$0.0197, $-$0.0024] & $-$0.0234, $-$0.0330, $-$0.0263 & no \\
\midrule
3.0 (secondary) & A: pre-trained $-$ random init. & $-$0.0283 [$-$0.0426, $-$0.0149] & $-$0.0295, $-$0.0345, $-$0.0427 & no \\
 & B: pre-trained $-$ ridge M3 & $-$0.0356 [$-$0.0509, $-$0.0215] & $-$0.0458, $-$0.0550, $-$0.0637 & no \\
\bottomrule
\end{tabularx}
% Internal provenance (not printed): pretrain-p3-2026-10-22/results.json
\endgroup

\end{table}


\begin{table}[htbp]
  \centering
  \caption{Macro AUROC with and without pre-training on tile labels}\label{tab:p3auroc}
  % Generated by build_ch4_assets.py - do not edit.
\begingroup\small
\renewcommand{\arraystretch}{1.15}
\begin{tabularx}{\textwidth}{>{\raggedright\arraybackslash}Xrrr}
\toprule
Training configuration & Ridge M3 & Network, random init. & Network, pre-trained \\
\midrule
Relaxed threshold 4.0 (main) & 0.8208 & 0.8175 & 0.8098 \\
Threshold 3.0 (secondary) & 0.8082 & 0.8009 & 0.7726 \\
\bottomrule
\end{tabularx}
\par\smallskip{\footnotesize Source: \nolinkurl{pretrain-p3-2026-10-22/results.json}.}
\endgroup

\end{table}



\section{Scope Comparison with MetaDegron}
MetaDegron \citep{zheng2024} was audited as a possible competitor and label source. Its
labels are binary motif-level annotations of curated degrons, not $\Delta$PSI; the audit therefore compares scope
and data, not predictive scores on the GPS matrix, because MetaDegron-D could not be run on GPS peptides
(Table~\ref{tab:metadegron}). Among the main findings, the curated positive set has only \nMdHuman{} human rows, and
almost no region of it coincides with the GPS peptides used here (\nMdOvScan{} scanning peptides, \nMdOvSat{} saturation peptide
and \nMdOvCrispr{} from the CRISPR screens). Only \nMdEThreeStrict{} E3 motif group keeps at least five samples after the strictest filter, so ``21 E3
ligases'' is a list of motif-compatible assignments, not 21 independently labelled classes. The reported
MetaDegron-X AUC of \nMdAuc{} was recomputed from the stored trees, but only on the authors' prepared tables, and \nMdInProps{} of the \nMdTestPos{}
positive test coordinates are in the property table of the fitting data; it is not an independent result and not a
GPS result. A simple composition logistic baseline on the same positives, split by protein cluster at cov80,
gives AUROC between \nCompLo{} and \nCompHi{} (core composition) and between \nCtxLo{} and \nCtxHi{} (29-residue
context), on a reconstructed population and not on the authors' own training set. The audit therefore suggests positioning MetaDegron as a complementary tool for motif and E3 information rather than a
head-to-head competitor on residue-level effect prediction; binary motif labels are not converted into $\Delta$PSI.

\begin{table}[htbp]
  \centering
  \caption[Audit of the MetaDegron data and overlap with the data of this thesis]{Audit of the MetaDegron data and overlap with the data of this thesis}\label{tab:metadegron}
  % Generated by build_ch4_assets.py - do not edit.
\begingroup\small
\renewcommand{\arraystretch}{1.15}
\begin{tabularx}{\textwidth}{>{\raggedright\arraybackslash}Xr}
\toprule
Audit item & Value \\
\midrule
Curated positive rows / distinct accessions & 302 / 240 \\
Rows with species suffix HUMAN & 227 \\
Rows in the property table used by MetaDegron-X & 212 \\
E3 motif groups with $\ge5$ samples: before filtering / after the strictest filter (human only, no regional or cluster overlap) & 13 / 1 \\
Same-region overlap with GPS scanning peptides (2\,624 peptides) & 2 \\
Same-region overlap with saturation peptides (248 mapped) & 1 \\
Same-region overlap with CRISPR-screen peptides (31 peptides) & 0 \\
MetaDegron-X AUC recomputed from the stored trees on the authors' prepared tables (not a GPS result) & 0.9670 \\
Positive test coordinates also present in the property table of the fitting data & 8 / 20 \\
\bottomrule
\end{tabularx}
\par\smallskip{\footnotesize Source: \nolinkurl{metadegron-analysis-2026-10-21/dataset-audit.json}, \nolinkurl{metadegron-analysis-2026-10-21/overlap-audit.json}, \nolinkurl{metadegron-analysis-2026-10-21/safe-x-reproduction.json}.}
\endgroup

\end{table}

\section{Additional Tables and Figures}
The following tables preserve the complete campaign record, including controls that do not improve
the selected metrics. All W rows share the same 235 peptides, 221 groups and 122,412 cells. Each
three-seed row averages predictions per cell before calculating metrics. The two Kaggle frozen-head
tuning controls have only seed 17; their scores are not directly equivalent to three-seed ensembles.

Training arms A, B and C restrict WT PSI to at most 3, 4 and 5, respectively; D uses the full range
and E balances the full range by stratum and peptide. Validation/test remain full-range. In the family
grid, ``only'' means pLM features with the common auxiliary features, and ``context'' additionally
includes ordered W5 features. Pointwise, CNN and multi-window CNN entries use frozen ESM2 features
on Kaggle; the local ``frozen'' and ``LoRA'' entries use the matched neural head described in
Section~\ref{sec:campaign-adaptation}. ``Rep1'' and ``mean12'' identify scanning-label pre-training
targets. Nested-family rows evaluate inner selection, not a retrospective outer-score winner.

For ridge controls, ``pair'' encodes WT/replacement identity in separate 20-category blocks;
it is not an interaction model and has no positional context, despite the saved W5 suffix.
``Composition'' adds whole-peptide amino acid frequencies, ``unordered'' adds local frequencies,
and ``ordered'' adds position-specific context. ``Sat-quality'' changes training weights using
replicate disagreement. ``Scan-rep1'' and ``scan-mean12'' add source-weighted scanning examples
to ridge training with an explicit source indicator; the CNN-transfer rows instead use scanning
pre-training. These are distinct interventions and should not be pooled into one transfer claim.

% Generated by build_campaign_assets.py; source: research/output-review-2026-11-06/results.json
% SHA256: ad5ba4a4bbffac32e090ca490893b728d9c4239f35484f9b6c6db13c76e94872

\begingroup\small\setlength{\tabcolsep}{3pt}
\begin{longtable}{@{}P{1.35cm}P{5.3cm}rrrrr@{}}
\caption{All saved OOF configurations. $s$ is the number of seeds; MAE and event ranking are peptide-macro. Only rows with $s=3$ are three-seed ensembles.}\label{tab:campaign-all}\\
\toprule Origin & Recipe & $s$ & MAE & AUC$+$ & AP$+$ & AUC$-$ \\ \midrule\endfirsthead
\toprule Origin & Recipe & $s$ & MAE & AUC$+$ & AP$+$ & AUC$-$ \\ \midrule\endhead
\midrule\multicolumn{7}{r}{Continued on next page}\endfoot
\bottomrule\endlastfoot
Kaggle & cnn-\allowbreak{}transfer-\allowbreak{}rep1-\allowbreak{}D & 3 & 0.3872 & 0.8425 & 0.5024 & 0.7132 \\
Kaggle & cnn-\allowbreak{}transfer-\allowbreak{}mean12-\allowbreak{}D & 3 & 0.3873 & 0.8391 & 0.4985 & 0.7185 \\
Kaggle & cnn-\allowbreak{}transfer-\allowbreak{}scratch-\allowbreak{}D & 3 & 0.3936 & 0.8291 & 0.4597 & 0.7098 \\
Kaggle & composition-\allowbreak{}w5-\allowbreak{}B & 3 & 0.4342 & 0.8111 & 0.4040 & 0.7456 \\
Kaggle & composition-\allowbreak{}w5-\allowbreak{}D & 3 & 0.4075 & 0.8109 & 0.4047 & 0.7506 \\
Kaggle & ESM2-\allowbreak{}35M-\allowbreak{}cnn-\allowbreak{}B & 3 & 0.3880 & 0.8506 & 0.5097 & 0.7484 \\
Kaggle & ESM2-\allowbreak{}35M-\allowbreak{}cnn-\allowbreak{}D & 3 & 0.3844 & 0.8488 & 0.5095 & 0.7550 \\
Kaggle & ESM2-\allowbreak{}35M-\allowbreak{}cnn-\allowbreak{}E & 3 & 0.3856 & 0.8474 & 0.5066 & 0.7540 \\
Kaggle & ESM2-\allowbreak{}35M-\allowbreak{}context-\allowbreak{}B & 3 & 0.4281 & 0.8432 & 0.4924 & 0.7095 \\
Kaggle & ESM2-\allowbreak{}35M-\allowbreak{}context-\allowbreak{}D & 3 & 0.4023 & 0.8445 & 0.4954 & 0.7301 \\
Kaggle & ESM2-\allowbreak{}35M-\allowbreak{}frozen-\allowbreak{}B & 1 & 0.4071 & 0.8233 & 0.4484 & 0.7512 \\
Kaggle & ESM2-\allowbreak{}35M-\allowbreak{}frozen-\allowbreak{}D & 1 & 0.3919 & 0.8221 & 0.4486 & 0.7657 \\
Kaggle & ESM2-\allowbreak{}35M-\allowbreak{}mcnn-\allowbreak{}B & 3 & 0.3855 & 0.8499 & 0.5036 & 0.7469 \\
Kaggle & ESM2-\allowbreak{}35M-\allowbreak{}mcnn-\allowbreak{}D & 3 & 0.3850 & 0.8463 & 0.5042 & 0.7570 \\
Kaggle & ESM2-\allowbreak{}35M-\allowbreak{}mcnn-\allowbreak{}E & 3 & 0.3838 & 0.8470 & 0.5064 & 0.7591 \\
Kaggle & ESM2-\allowbreak{}35M-\allowbreak{}only-\allowbreak{}B & 3 & 0.4158 & 0.8109 & 0.4291 & 0.7827 \\
Kaggle & ESM2-\allowbreak{}35M-\allowbreak{}only-\allowbreak{}D & 3 & 0.4008 & 0.8131 & 0.4340 & 0.7868 \\
Kaggle & ESM2-\allowbreak{}35M-\allowbreak{}pointwise-\allowbreak{}B & 3 & 0.4004 & 0.8375 & 0.4656 & 0.7440 \\
Kaggle & ESM2-\allowbreak{}35M-\allowbreak{}pointwise-\allowbreak{}D & 3 & 0.3897 & 0.8375 & 0.4739 & 0.7515 \\
Kaggle & ESM2-\allowbreak{}35M-\allowbreak{}pointwise-\allowbreak{}E & 3 & 0.3886 & 0.8357 & 0.4712 & 0.7557 \\
Kaggle & Nested family & 3 & 0.3975 & 0.8367 & 0.4788 & 0.7578 \\
Kaggle & onehot-\allowbreak{}cnn-\allowbreak{}B & 3 & 0.3997 & 0.8304 & 0.4575 & 0.7030 \\
Kaggle & onehot-\allowbreak{}cnn-\allowbreak{}D & 3 & 0.3936 & 0.8291 & 0.4597 & 0.7098 \\
Kaggle & onehot-\allowbreak{}cnn-\allowbreak{}E & 3 & 0.3947 & 0.8267 & 0.4564 & 0.7166 \\
Kaggle & onehot-\allowbreak{}mcnn-\allowbreak{}B & 3 & 0.3985 & 0.8267 & 0.4516 & 0.7046 \\
Kaggle & onehot-\allowbreak{}mcnn-\allowbreak{}D & 3 & 0.3935 & 0.8293 & 0.4588 & 0.7199 \\
Kaggle & onehot-\allowbreak{}mcnn-\allowbreak{}E & 3 & 0.3921 & 0.8289 & 0.4592 & 0.7200 \\
Kaggle & onehot-\allowbreak{}pointwise-\allowbreak{}B & 3 & 0.4059 & 0.8245 & 0.4299 & 0.7129 \\
Kaggle & onehot-\allowbreak{}pointwise-\allowbreak{}D & 3 & 0.3955 & 0.8269 & 0.4424 & 0.7264 \\
Kaggle & onehot-\allowbreak{}pointwise-\allowbreak{}E & 3 & 0.3920 & 0.8235 & 0.4426 & 0.7338 \\
Kaggle & ordered-\allowbreak{}w3-\allowbreak{}B & 3 & 0.4331 & 0.8280 & 0.4511 & 0.6927 \\
Kaggle & ordered-\allowbreak{}w5-\allowbreak{}A & 3 & 0.4711 & 0.8192 & 0.4436 & 0.6602 \\
Kaggle & ordered-\allowbreak{}w5-\allowbreak{}B & 3 & 0.4336 & 0.8312 & 0.4632 & 0.6745 \\
Kaggle & ordered-\allowbreak{}w5-\allowbreak{}C & 3 & 0.4177 & 0.8314 & 0.4606 & 0.6875 \\
Kaggle & ordered-\allowbreak{}w5-\allowbreak{}D & 3 & 0.4078 & 0.8319 & 0.4614 & 0.6931 \\
Kaggle & ordered-\allowbreak{}w5-\allowbreak{}D-\allowbreak{}fraction0.25 & 3 & 0.4154 & 0.8231 & 0.4468 & 0.6925 \\
Kaggle & ordered-\allowbreak{}w5-\allowbreak{}D-\allowbreak{}fraction0.5 & 3 & 0.4113 & 0.8280 & 0.4525 & 0.6975 \\
Kaggle & ordered-\allowbreak{}w5-\allowbreak{}D-\allowbreak{}fraction0.75 & 3 & 0.4083 & 0.8297 & 0.4555 & 0.6983 \\
Kaggle & ordered-\allowbreak{}w5-\allowbreak{}E & 3 & 0.4036 & 0.8275 & 0.4499 & 0.7109 \\
Kaggle & pair-\allowbreak{}w5-\allowbreak{}B & 3 & 0.4149 & 0.8093 & 0.4024 & 0.7458 \\
Kaggle & pair-\allowbreak{}w5-\allowbreak{}D & 3 & 0.4017 & 0.8092 & 0.4040 & 0.7499 \\
Kaggle & sat-\allowbreak{}only-\allowbreak{}D & 3 & 0.4078 & 0.8319 & 0.4614 & 0.6931 \\
Kaggle & sat-\allowbreak{}quality-\allowbreak{}D & 3 & 0.4118 & 0.8307 & 0.4579 & 0.6918 \\
Kaggle & scan-\allowbreak{}mean12-\allowbreak{}D & 3 & 0.4048 & 0.8334 & 0.4643 & 0.6963 \\
Kaggle & scan-\allowbreak{}rep1-\allowbreak{}D & 3 & 0.4047 & 0.8334 & 0.4643 & 0.6964 \\
Kaggle & unordered-\allowbreak{}w5-\allowbreak{}B & 3 & 0.4316 & 0.8260 & 0.4414 & 0.7096 \\
Kaggle & unordered-\allowbreak{}w5-\allowbreak{}D & 3 & 0.4060 & 0.8258 & 0.4399 & 0.7199 \\
Local & Ankh-\allowbreak{}base-\allowbreak{}context-\allowbreak{}B & 3 & 0.4279 & 0.8426 & 0.4949 & 0.7097 \\
Local & Ankh-\allowbreak{}base-\allowbreak{}context-\allowbreak{}D & 3 & 0.4035 & 0.8440 & 0.4988 & 0.7310 \\
Local & Ankh-\allowbreak{}base-\allowbreak{}only-\allowbreak{}B & 3 & 0.4188 & 0.8075 & 0.4350 & 0.7822 \\
Local & Ankh-\allowbreak{}base-\allowbreak{}only-\allowbreak{}D & 3 & 0.4048 & 0.8088 & 0.4396 & 0.7882 \\
Local & ESM1b-\allowbreak{}650M-\allowbreak{}context-\allowbreak{}B & 3 & 0.4295 & 0.8359 & 0.4773 & 0.6814 \\
Local & ESM1b-\allowbreak{}650M-\allowbreak{}context-\allowbreak{}D & 3 & 0.4069 & 0.8360 & 0.4752 & 0.7032 \\
Local & ESM1b-\allowbreak{}650M-\allowbreak{}only-\allowbreak{}B & 3 & 0.4203 & 0.8016 & 0.4005 & 0.7431 \\
Local & ESM1b-\allowbreak{}650M-\allowbreak{}only-\allowbreak{}D & 3 & 0.4054 & 0.8024 & 0.4088 & 0.7547 \\
Local & ESM1v-\allowbreak{}650M-\allowbreak{}context-\allowbreak{}B & 3 & 0.4289 & 0.8384 & 0.4838 & 0.7009 \\
Local & ESM1v-\allowbreak{}650M-\allowbreak{}context-\allowbreak{}D & 3 & 0.4050 & 0.8378 & 0.4814 & 0.7187 \\
Local & ESM1v-\allowbreak{}650M-\allowbreak{}only-\allowbreak{}B & 3 & 0.4179 & 0.8022 & 0.4153 & 0.7677 \\
Local & ESM1v-\allowbreak{}650M-\allowbreak{}only-\allowbreak{}D & 3 & 0.4032 & 0.8027 & 0.4194 & 0.7725 \\
Local & ESM2-\allowbreak{}150M-\allowbreak{}context-\allowbreak{}B & 3 & 0.4296 & 0.8381 & 0.4894 & 0.7024 \\
Local & ESM2-\allowbreak{}150M-\allowbreak{}context-\allowbreak{}D & 3 & 0.4041 & 0.8390 & 0.4856 & 0.7199 \\
Local & ESM2-\allowbreak{}150M-\allowbreak{}only-\allowbreak{}B & 3 & 0.4183 & 0.8029 & 0.4229 & 0.7677 \\
Local & ESM2-\allowbreak{}150M-\allowbreak{}only-\allowbreak{}D & 3 & 0.4031 & 0.8034 & 0.4258 & 0.7729 \\
Local & ESM2-\allowbreak{}35M-\allowbreak{}context-\allowbreak{}B & 3 & 0.4280 & 0.8432 & 0.4921 & 0.7097 \\
Local & ESM2-\allowbreak{}35M-\allowbreak{}context-\allowbreak{}D & 3 & 0.4022 & 0.8445 & 0.4954 & 0.7302 \\
Local & ESM2-\allowbreak{}35M-\allowbreak{}frozen-\allowbreak{}B & 3 & 0.4034 & 0.8263 & 0.4554 & 0.7541 \\
Local & ESM2-\allowbreak{}35M-\allowbreak{}frozen-\allowbreak{}D & 3 & 0.3904 & 0.8246 & 0.4544 & 0.7667 \\
Local & ESM2-\allowbreak{}35M-\allowbreak{}lora-\allowbreak{}B & 3 & 0.3637 & 0.8550 & 0.5517 & 0.7921 \\
Local & ESM2-\allowbreak{}35M-\allowbreak{}lora-\allowbreak{}D & 3 & 0.3650 & 0.8536 & 0.5474 & 0.8028 \\
Local & ESM2-\allowbreak{}35M-\allowbreak{}only-\allowbreak{}B & 3 & 0.4157 & 0.8109 & 0.4290 & 0.7829 \\
Local & ESM2-\allowbreak{}35M-\allowbreak{}only-\allowbreak{}D & 3 & 0.4007 & 0.8131 & 0.4340 & 0.7870 \\
Local & ESM2-\allowbreak{}650M-\allowbreak{}context-\allowbreak{}B & 3 & 0.4319 & 0.8392 & 0.4832 & 0.6923 \\
Local & ESM2-\allowbreak{}650M-\allowbreak{}context-\allowbreak{}D & 3 & 0.4071 & 0.8394 & 0.4796 & 0.7139 \\
Local & ESM2-\allowbreak{}650M-\allowbreak{}only-\allowbreak{}B & 3 & 0.4204 & 0.8033 & 0.4105 & 0.7663 \\
Local & ESM2-\allowbreak{}650M-\allowbreak{}only-\allowbreak{}D & 3 & 0.4056 & 0.8022 & 0.4093 & 0.7707 \\
Local & ESM2-\allowbreak{}8M-\allowbreak{}context-\allowbreak{}B & 3 & 0.4288 & 0.8419 & 0.4943 & 0.7052 \\
Local & ESM2-\allowbreak{}8M-\allowbreak{}context-\allowbreak{}D & 3 & 0.4018 & 0.8433 & 0.4948 & 0.7215 \\
Local & ESM2-\allowbreak{}8M-\allowbreak{}only-\allowbreak{}B & 3 & 0.4150 & 0.8107 & 0.4398 & 0.7799 \\
Local & ESM2-\allowbreak{}8M-\allowbreak{}only-\allowbreak{}D & 3 & 0.3995 & 0.8128 & 0.4375 & 0.7839 \\
Local & ESM-\allowbreak{}C-\allowbreak{}300M-\allowbreak{}context-\allowbreak{}B & 3 & 0.4294 & 0.8379 & 0.4806 & 0.6937 \\
Local & ESM-\allowbreak{}C-\allowbreak{}300M-\allowbreak{}context-\allowbreak{}D & 3 & 0.4038 & 0.8391 & 0.4820 & 0.7201 \\
Local & ESM-\allowbreak{}C-\allowbreak{}300M-\allowbreak{}only-\allowbreak{}B & 3 & 0.4203 & 0.8001 & 0.4038 & 0.7643 \\
Local & ESM-\allowbreak{}C-\allowbreak{}300M-\allowbreak{}only-\allowbreak{}D & 3 & 0.4046 & 0.8011 & 0.4062 & 0.7750 \\
Local & Nested family & 3 & 0.3945 & 0.8409 & 0.4877 & 0.7530 \\
Local & ProtT5-\allowbreak{}context-\allowbreak{}B & 3 & 0.4301 & 0.8414 & 0.4874 & 0.7011 \\
Local & ProtT5-\allowbreak{}context-\allowbreak{}D & 3 & 0.4055 & 0.8425 & 0.4900 & 0.7244 \\
Local & ProtT5-\allowbreak{}only-\allowbreak{}B & 3 & 0.4211 & 0.8036 & 0.4143 & 0.7647 \\
Local & ProtT5-\allowbreak{}only-\allowbreak{}D & 3 & 0.4056 & 0.8049 & 0.4223 & 0.7740 \\
\end{longtable}\endgroup


\Needspace{11\baselineskip}

The following supporting tables were moved here from Chapter~\ref{chap:evaluation}; their numbers,
denominators and intervals are unchanged: the newly added groups of the learning-curve study
(Table~\ref{tab:newgroup}), the robustness of the threshold sensitivity (Table~\ref{tab:partcrobust}) and
the residual breakdowns of Section~\ref{sec:residual-diagnostics} (Tables~\ref{tab:progress-effect-errors}
and~\ref{tab:progress-largest-errors}).

\begin{table}[htbp]
  \centering
  \caption[Cells of the benchmark peptides and of the 85 peptides added by relaxing the threshold]{Cells of the benchmark peptides and of the 85 peptides added by relaxing the threshold}\label{tab:newgroup}
  % Generated by build_ch4_assets.py - do not edit.
\begingroup\small
\renewcommand{\arraystretch}{1.15}
\begin{tabularx}{\textwidth}{>{\raggedright\arraybackslash}Xrr}
\toprule
Statistic & Benchmark set (91 peptides) & Added set (85 peptides) \\
\midrule
Median $\mathrm{PSI}_{\mathrm{WT}}$ & 2.593 & 3.346 \\
Fraction of positive cells ($\Delta\mathrm{PSI}\ge 1$) & 0.1310 & 0.1140 \\
Median $\Delta$PSI & 0.125 & 0.083 \\
95th percentile of $\Delta$PSI & 1.697 & 1.507 \\
\bottomrule
\end{tabularx}
% Internal provenance (not printed): learning-curve-2026-10-22/results.json
\endgroup

\end{table}

\begin{table}[htbp]
  \centering
  \caption{Robustness of the gain at threshold 4.0 (ridge M3)}\label{tab:partcrobust}
  % Generated by build_ch4_assets.py - do not edit.
\begingroup\small
\renewcommand{\arraystretch}{1.15}
\begin{tabularx}{\textwidth}{>{\raggedright\arraybackslash}Xrr}
\toprule
Condition (threshold 4.0) & Ridge $\Delta$ vs. base [95\% CI] & CI $>0$ \\
\midrule
cov80, $\Delta\mathrm{PSI}\ge 0.5$ & +0.0110 [+0.0051, +0.0169] & yes \\
cov80, $\Delta\mathrm{PSI}\ge 1.0$ (main) & +0.0126 [+0.0061, +0.0190] & yes \\
cov80, $\Delta\mathrm{PSI}\ge 1.5$ & +0.0094 [+0.0037, +0.0149] & yes \\
cov50 fold assignment, $\Delta\mathrm{PSI}\ge 1.0$ & +0.0096 [+0.0032, +0.0159] & yes \\
\bottomrule
\end{tabularx}
\par\smallskip{\footnotesize Source: \nolinkurl{learning-curve-2026-10-22/results-c.json}.}
\endgroup

\end{table}

\begin{table}[htbp]
  \centering
  \caption[OOF MAE conditional on the measured effect class ($\Delta=\Delta$PSI)]{OOF MAE conditional on the measured effect class ($\Delta=\Delta$PSI).
  Each metric averages errors within each eligible peptide, then over peptides in that class.
  Support counts differ from AUROC/average-precision eligibility, which requires at least three
  cells of each class; the fixed-budget precision replay retains all peptides.}
  \label{tab:progress-effect-errors}
  % Generated by build_progress_assets.py; descriptive OOF diagnostics.
% Source: research/progress-review-2026-11-06/local-audit.json
% SHA256: 2ea1d4235a4123bdcd0411f788f1356be6fe99b835105d1743700cf2670195a5
{\small\setlength{\tabcolsep}{3.5pt}
\begin{tabularx}{\textwidth}{Xrrrrrrr}
\toprule
Effect & Cells & Pep./grp. & Fr.-B & LoRA-B & Fr.-D & LoRA-D & Zero \\
\midrule
$-1<\Delta<1$ & 108,782 & 235/221 & 0.3270 & 0.2879 & 0.3045 & 0.2843 & 0.3031 \\
$\Delta\ge1$ & 12,053 & 198/190 & 0.8187 & 0.7997 & 0.8972 & 0.8479 & 1.4155 \\
$\Delta\le-1$ & 1,577 & 151/144 & 1.2407 & 1.0385 & 1.1394 & 0.9510 & 1.2320 \\
\bottomrule
\end{tabularx}}

\end{table}

\begin{table}[htbp]
  \centering
  \caption[Three largest absolute cell errors for LoRA-B on W, selected after inspection]{Three largest absolute cell errors for LoRA-B on W, selected after inspection.
  Positions are within the measured peptide; measured and predicted values are $\Delta$PSI.
  These illustrate failures and are not clinical or mechanistic case validations.}
  \label{tab:progress-largest-errors}
  % Generated by build_progress_assets.py; descriptive OOF diagnostics.
% Source: research/progress-review-2026-11-06/local-audit.json
% SHA256: 2ea1d4235a4123bdcd0411f788f1356be6fe99b835105d1743700cf2670195a5
{\small\setlength{\tabcolsep}{4pt}
\begin{tabularx}{\textwidth}{Xrrrrr}
\toprule
Peptide & Fold & Position & Change & Measured & Predicted \\
\midrule
P32881\_IFNA8\_5 & 2 & 17 & Y$\to$F & 3.2866 & -0.1573 \\
P13942\_COL11A2\_20 & 1 & 18 & R$\to$I & 3.1698 & -0.2479 \\
P32881\_IFNA8\_5 & 2 & 17 & Y$\to$M & 3.0830 & -0.2058 \\
\bottomrule
\end{tabularx}}

\end{table}

\FloatBarrier
\section{Seed Stability and Inner Selection}
\begin{table}[htbp]
  \centering
  \caption[Single-seed and ensemble peptide-macro MAE for the matched local core experiment]{Single-seed and ensemble peptide-macro MAE for the matched local core experiment.}
  \label{tab:campaign-seeds}
  % Generated by build_campaign_assets.py; source: research/output-review-2026-11-06/results.json
% SHA256: ad5ba4a4bbffac32e090ca490893b728d9c4239f35484f9b6c6db13c76e94872

{\small\setlength{\tabcolsep}{4pt}
\begin{tabularx}{\textwidth}{Xrrrr}
\toprule
Recipe & Seed 17 & Seed 29 & Seed 43 & Ensemble \\
\midrule
Frozen-B & 0.4073 & 0.4031 & 0.4059 & 0.4034 \\
Frozen-D & 0.3920 & 0.3918 & 0.3938 & 0.3904 \\
LoRA-B & 0.3794 & 0.3786 & 0.3779 & 0.3637 \\
LoRA-D & 0.3794 & 0.3778 & 0.3788 & 0.3650 \\
\bottomrule
\end{tabularx}}

\end{table}
\begin{table}[htbp]
  \centering
  \caption[Inner-selected epoch counts over 15 outer fits per local core recipe]{Inner-selected epoch counts over 15 outer fits per local core recipe.}
  \label{tab:campaign-epochs}
  % Generated by build_campaign_assets.py; source: research/output-review-2026-11-06/results.json
% SHA256: ad5ba4a4bbffac32e090ca490893b728d9c4239f35484f9b6c6db13c76e94872

{\small\setlength{\tabcolsep}{4pt}
\begin{tabularx}{\textwidth}{Xrrrr}
\toprule
Recipe & Folds & Minimum & Median & Maximum \\
\midrule
frozen-B & 15 & 1 & 1 & 3 \\
frozen-D & 15 & 1 & 1 & 3 \\
lora-B & 15 & 1 & 3 & 7 \\
lora-D & 15 & 1 & 2 & 4 \\
\bottomrule
\end{tabularx}}

\end{table}
\begin{table}[htbp]
  \centering
  \caption[Local frozen-family recipes chosen by the saved inner-MSE selector]{Local frozen-family recipes chosen by the saved inner-MSE selector.
  Frequencies across seed/fold pairs are not independent votes for a model family.}
  \label{tab:campaign-nested}
  % Generated by build_campaign_assets.py; source: research/output-review-2026-11-06/results.json
% SHA256: ad5ba4a4bbffac32e090ca490893b728d9c4239f35484f9b6c6db13c76e94872

{\small\setlength{\tabcolsep}{4pt}
\begin{tabularx}{\textwidth}{Xr}
\toprule
Selected recipe & Outer fits (of 15) \\
\midrule
ESM2-\allowbreak{}8M-\allowbreak{}context-\allowbreak{}D & 3 \\
Ankh-\allowbreak{}base-\allowbreak{}context-\allowbreak{}D & 2 \\
ESM2-\allowbreak{}35M-\allowbreak{}context-\allowbreak{}D & 2 \\
ESM2-\allowbreak{}8M-\allowbreak{}only-\allowbreak{}D & 2 \\
Ankh-\allowbreak{}base-\allowbreak{}only-\allowbreak{}D & 1 \\
ESM1b-\allowbreak{}650M-\allowbreak{}context-\allowbreak{}D & 1 \\
ESM1v-\allowbreak{}650M-\allowbreak{}only-\allowbreak{}D & 1 \\
ESM2-\allowbreak{}150M-\allowbreak{}only-\allowbreak{}D & 1 \\
ESM2-\allowbreak{}35M-\allowbreak{}only-\allowbreak{}D & 1 \\
ESM-\allowbreak{}C-\allowbreak{}300M-\allowbreak{}context-\allowbreak{}B & 1 \\
\bottomrule
\end{tabularx}}

\end{table}

\FloatBarrier
\section{Retrospective Prioritisation Sensitivity}
The following tables supplement Section~\ref{sec:res-prioritization}. Precision includes every
W peptide, including those without a positive hit. Seeds share outer folds; the ensemble averages
predictions before ranking. Neither a seed nor a budget is selected after inspecting utility.
The 100 fixed tie seeds give unchanged results because no exact score tie crosses a tested
cutoff. Checks retain all no-hit cases and independently reconstruct all 14,100 ensemble
peptide--policy--endpoint--budget rows. They verify arithmetic and saved evidence, not new
biological data or the complete historical model-selection process.

\begin{table}[htbp]
  \centering
  \caption[Positive Precision@10 in percent for each seed and the three-seed ensemble]{Positive Precision@10 in percent for each seed and the three-seed ensemble.
  Every row and column uses 2,350 selections on 235 peptides.}
  \label{tab:priority-seeds}
  % Generated by build_prioritization_assets.py; retrospective development OOF.
% Source: research/oof-prioritization-2026-11-07/results.json
% SHA256: 4016cb5686f3520956cdc49268e2ea5363e678a6f1ce8555cec53e19a6938b87
{\small\setlength{\tabcolsep}{4pt}
\begin{tabularx}{\textwidth}{Xrrrr}
\toprule
Policy & Seed 17 & Seed 29 & Seed 43 & Ensemble \\
\midrule
LoRA-D & 51.91 & 50.13 & 52.77 & 53.62 \\
Ridge W5-D & 46.26 & 46.26 & 46.26 & 46.26 \\
Frozen-D & 46.98 & 47.70 & 47.66 & 47.83 \\
LoRA-B & 51.32 & 51.62 & 51.19 & 53.96 \\
\bottomrule
\end{tabularx}}

\end{table}
\begin{table}[htbp]
  \centering
  \caption[Positive peptide-macro precision in percent across the six fixed analysis budgets]{Positive peptide-macro precision in percent across the six fixed analysis budgets.
  Total selections are equal across policies within each row. Fractional budgets use ceiling
  rounding within each peptide and differ from a fixed ten-mutation budget.}
  \label{tab:priority-budgets}
  % Generated by build_prioritization_assets.py; retrospective development OOF.
% Source: research/oof-prioritization-2026-11-07/results.json
% SHA256: 4016cb5686f3520956cdc49268e2ea5363e678a6f1ce8555cec53e19a6938b87
{\footnotesize\setlength{\tabcolsep}{4pt}
\begin{tabularx}{\textwidth}{Xrrrrrr}
\toprule
Budget & Total & LoRA-D & W5-D & Random & Frozen-D & LoRA-B \\
\midrule
$k=5$ & 1,175 & 54.89 & 47.91 & 9.78 & 49.02 & 55.66 \\
$k=10$ & 2,350 & 53.62 & 46.26 & 9.78 & 47.83 & 53.96 \\
$k=20$ & 4,700 & 50.55 & 43.04 & 9.78 & 43.55 & 51.06 \\
Top 5\% & 6,225 & 49.04 & 40.80 & 9.78 & 40.81 & 48.95 \\
Top 10\% & 12,399 & 41.50 & 34.79 & 9.78 & 34.76 & 41.59 \\
Top 20\% & 24,611 & 30.32 & 27.49 & 9.78 & 27.75 & 30.61 \\
\bottomrule
\end{tabularx}}

\end{table}
\FloatBarrier
\Needspace{11\baselineskip}

Tables~\ref{tab:priority-recall-budget} and~\ref{tab:priority-strata} were moved here from
Section~\ref{sec:res-prioritization} to keep the results chapter readable; their numbers are unchanged.

\begin{table}[htbp]
  \centering
  \caption[Mean selections to reach a known positive-hit recall, on 198 peptides with hits]{Mean selections to reach a known positive-hit recall, on 198 peptides with hits.
  Each peptide target is rounded upward. Random values are expected order-statistic ranks.}
  \label{tab:priority-recall-budget}
  % Generated by build_prioritization_assets.py; retrospective development OOF.
% Source: research/oof-prioritization-2026-11-07/results.json
% SHA256: 4016cb5686f3520956cdc49268e2ea5363e678a6f1ce8555cec53e19a6938b87
{\small\setlength{\tabcolsep}{4pt}
\begin{tabularx}{\textwidth}{Xrr}
\toprule
Policy & 50\% hit recall & 80\% hit recall \\
\midrule
LoRA-D & 71.32 & 154.36 \\
Ridge W5-D & 81.99 & 170.53 \\
Random expectation & 255.51 & 404.48 \\
\bottomrule
\end{tabularx}}

\end{table}

\begin{table}[htbp]
  \centering
  \caption[Positive Precision@10 by WT-PSI stratum, in percent]{Positive Precision@10 by WT-PSI stratum, in percent. Peptides without positive hits
  remain in precision; zero-event strata do not supply recall or enrichment evidence.}
  \label{tab:priority-strata}
  % Generated by build_prioritization_assets.py; retrospective development OOF.
% Source: research/oof-prioritization-2026-11-07/results.json
% SHA256: 4016cb5686f3520956cdc49268e2ea5363e678a6f1ce8555cec53e19a6938b87
{\small\setlength{\tabcolsep}{4pt}
\begin{tabularx}{\textwidth}{Xrrrr}
\toprule
WT PSI & Peptides & No (+) hit & LoRA-D P@10 & W5-D P@10 \\
\midrule
$\le3$ & 91 & 0 & 71.10 & 62.64 \\
$(3,4]$ & 85 & 2 & 63.76 & 52.59 \\
$(4,5]$ & 36 & 12 & 19.72 & 19.44 \\
$>5$ & 23 & 23 & 0.00 & 0.00 \\
\bottomrule
\end{tabularx}}

\end{table}

\FloatBarrier
\section{Execution Cost and Output Inventory}
\begin{table}[htbp]
  \centering
  \caption[Recorded durations of completed local job attempts, counting inherited work once]{Recorded durations of completed local job attempts, counting inherited work once.
  Earlier attempts of paused/resumed jobs may be absent; these are not GPU-active or full elapsed times.}
  \label{tab:campaign-timing}
  % Generated by build_campaign_assets.py; source: research/output-review-2026-11-06/results.json
% SHA256: ad5ba4a4bbffac32e090ca490893b728d9c4239f35484f9b6c6db13c76e94872

{\small\setlength{\tabcolsep}{4pt}
\begin{tabularx}{\textwidth}{Xrrr}
\toprule
Job category & Jobs & Recorded h & Median min \\
\midrule
other & 15 & 3.07 & 0.19 \\
probe & 9 & 0.01 & 0.09 \\
extract & 9 & 0.96 & 5.35 \\
ridge & 540 & 1.42 & 0.16 \\
tune-\allowbreak{}frozen & 30 & 14.07 & 27.16 \\
tune-\allowbreak{}lora & 30 & 36.01 & 73.30 \\
\bottomrule
\end{tabularx}}

\end{table}
The cache contains pinned weights and executable source snapshots as well as metrics, OOF arrays,
features, fitted models and recovery checkpoints. Physical file counts in Table~\ref{tab:campaign-inventory}
can include duplicate content, partial files and superseded attempts. They must not be added to
the checksum-verified artifact totals as though they were separate scientific observations.
The transport archive contains the same expanded namespace and file sizes as the output tree;
it packages Kaggle and local campaigns together with cached weights. It is an additional physical
copy, not another experimental run. The archive review verifies small-member CRCs but does not
recheck every large compressed payload.
Zero-filled JSON metadata in the reviewed expanded copy were recovered from the transport backup,
using saved artifact hashes or exact agreement with state records before a fresh integrity check.
The historical timer-failure local run contributes no scientific OOF results. The interrupted local
source remains required because the successor references its inherited artifacts.

\begin{table}[htbp]
  \centering
  \caption[Physical output inventory at review time]{Physical output inventory at review time. Weight caches and code snapshots are supporting
  assets, not evaluation datasets.}
  \label{tab:campaign-inventory}
  % Generated by build_campaign_assets.py; source: research/output-review-2026-11-06/results.json
% SHA256: ad5ba4a4bbffac32e090ca490893b728d9c4239f35484f9b6c6db13c76e94872

{\small\setlength{\tabcolsep}{4pt}
\begin{tabularx}{\textwidth}{Xrr}
\toprule
Output role & Physical files & GiB \\
\midrule
Pinned executable snapshot & 30 & 0.000 \\
Superseded timer-failure run & 18 & 0.000 \\
Earlier local source campaign & 2959 & 3.541 \\
Completed local successor & 2754 & 2.131 \\
Nine-checkpoint local weight cache & 63 & 14.110 \\
Kaggle reference campaign & 2990 & 2.859 \\
Reference ESM2-35M weight cache & 8 & 0.127 \\
Top-level cache/session metadata & 5 & 0.000 \\
Duplicate transport archive & 1 & 17.542 \\
\bottomrule
\end{tabularx}}

\end{table}
